\chapter{Plano de Trabalho e Cronograma}
\label{ch:plano-trabalho}

% Meta editorial: 3--5 páginas. Capítulo curto, verificável e prospectivo.

\section{Estado atual}
% Tabela/inventário datado: escrita, implementação, experimentos concluídos,
% em execução, incompletos, falhos e análises disponíveis. Não inferir status
% apenas pela existência de configs; confirmar logs/reports.

\section{Trabalho restante até a defesa}
% Organizar por entregável: auditoria texto-código (inclusive nLocal/nGlobal),
% experimentos principais restantes, análise estatística, figuras, redação,
% revisão com orientação e preparação da defesa. Avg-K/SBS e extensões não podem
% bloquear o núcleo mínimo de validação da ANT sem justificativa explícita.

\section{Cronograma mensal e marcos}
% Tabela mês x atividade, com legenda planejado/em andamento/concluído e marcos
% observáveis. TODO: preencher datas somente após confirmar calendário do
% programa e alinhar o conjunto mínimo de experimentos com a orientação.

\section{Riscos e contingências}
% GPU/tempo, variância, experimentos incompletos, equivalência nLocal/nGlobal,
% resultados inconclusivos e redução de escopo. Definir plano mínimo para
% responder à hipótese principal mesmo que investigações auxiliares falhem.

\section{Critérios de conclusão}
% Condições para encerrar experimentação, análise, redação, revisão e defesa.
